Autoregressive language models train on all tokens in parallel, yet must generate them one at a time.
Diffusion language models remove this limit by denoising many tokens in each forward pass \citep{nie2026large, ye2025dream7bdiffusionlarge, team2026diffusiongemma}.

How well this works depends on the sampler.
In the continuous domain, many improvements have been driven by improved sampling algorithms \citep{song2020denoising}.
The same is true for the language domain, where sampling order has been shown to have a dramatic effect on task performance \citep{kim2025train, ben2026accelerated, kim2026klass}.

Any diffusion language sampler makes two choices for each position in its denoising window.
The first is \emph{where} the next token comes from: the model's prediction, the token already at that position, or fresh noise.
The second is \emph{when} the position is committed to the KV cache, after which it can no longer change.
Existing samplers tie the two choices together.
Samplers for masked diffusion freeze a token the moment it is accepted \citep{nie2026large, ben2026accelerated}, so acceptance and commitment are the same event.
Block samplers sit at the other extreme: every token waits until the whole block has converged \citep{arriola2025blockdiffusioninterpolatingautoregressive, team2026diffusiongemma}.

We separate the two choices and study each extreme on DiffusionGemma \citep{team2026diffusiongemma}, a uniform diffusion model whose default sampler is the entropy-bound (EB) sampler of \citet{ben2026accelerated}.
Both extremes are costly.
\emph{Committing too early} removes self-correction: when accepted tokens are frozen, macro accuracy falls from 73\% to 20\% as the entropy bound is relaxed, while the unfrozen sampler stays near 70\%.
This self-correction is a property of the model rather than the sampler, as it persists when renoising is removed.
\emph{Committing too late} wastes compute: tokens converge in a largely left-to-right order, so a converged prefix sits idle in the window, and after ten denoising steps more than half of the slots are spent on tokens that are already final.

These two findings motivate the Rolling Entropy Sampler (RES).
RES uses the entropy bound for both choices: it accepts tokens as the EB sampler does, and it commits the accepted prefix of the window as soon as that prefix is at least $s$ tokens long.
Every position outside the committed prefix stays free to change.
The minimum shift $s$ sets how long a token stays in the window, and gives a second knob, beside the entropy bound, for moving along the speed-quality frontier.

In summary, our contributions are as follows:
\begin{itemize}
    \item We write samplers as a routing rule and a commitment rule (Section~\ref{sec:motivation}), which separates accepting a token from committing it.
    \item We show that DiffusionGemma self-corrects without renoising, and that freezing accepted tokens destroys its speed-quality frontier (Section~\ref{sec:too-early}).
    \item We show that the EB sampler converges left to right, leaving over half of the window idle late in each block (Section~\ref{sec:too-late}).
    \item We introduce RES (Section~\ref{sec:method}). Across seven benchmarks it raises tokens per forward pass by 32\% at an average cost of 1.9 points of accuracy, and it extends the frontier beyond what relaxing the entropy bound alone reaches (Section~\ref{sec:experiments}).
\end{itemize}

% Frontier diffusion language models decode block-autoregressively: the model iteratively denoises a fixed canvas of tokens, and only once the entire canvas has converged is it committed and the next canvas begun \citep{arriola2025blockdiffusioninterpolatingautoregressive, team2026diffusiongemma, song2025seed, nie2026large, ye2025dream7bdiffusionlarge}. 
% Within a canvas, adaptive samplers such as the Entropy-Bounded (EB) sampler decide how many tokens to resolve per denoising step, concentrating work where the model is uncertain \citep{ben2026accelerated}. 
% The block structure, however, imposes a cost that no within-canvas sampler can avoid: every denoising step is a full-width forward pass, regardless of how little unresolved work the canvas still contains. 
% As a canvas nears convergence, the sampler processes the entire sequence, regardless of whether most of the positions are already resolved. 
% This property of block denoising represents a major inefficiency, as the compute expended on denoised positions could be redirected towards resolving noisy tokens in the next canvas. 

% Two observations suggest that this inefficiency can be addressed. First, the EB sampler commits tokens in a loosely autoregressive order: across [benchmarks], the order in which positions realize their final token, or when they are committed, tend to correlate with an autoregressive ordering. 
% Second, while uniform diffusion models can theoretically change the token identity even after being accepted during the EB sampler, we find that this tends to occur with very low probability. 

% These two observations motivate \algname. 
% Rather than viewing the denoising window as moving across the sequence, we view the sequence as flowing through a fixed denoising window.
% This behavior is governed by two rules: the standard acceptance rule introduced by the EB sampler, and a commitment rule: when the accepted positions include a contiguous prefix of the window, the window appends them to the KV cache and moves past them. The freed positions are refilled with fresh noise at the trailing edge of the window, so subsequent passes denoise unresolved tokens. 
% The fixed window also lets us recast the usual time-dependent temperature schedule as a static per-position profile: each window slot carries a fixed temperature, highest at the intake and lowest at the exit, so a token flowing through the window anneals from exploratory sampling to sharp predictions by construction, with no per-step schedule logic.

% Throughout this paper, we use DiffusionGemma~\citep{team2026diffusiongemma} as our primary model, though \algname applies to block diffusion language models generally. We begin with three empirical observations that motivate the design (Section~\ref{sec:motivation}). First, DiffusionGemma commits tokens in a strongly autoregressive order, consistent with observations across other diffusion LMs \citep{ni2026flexibility}. Second, tokens accepted by the entropy rule are rarely renoised in later steps: the model retains high confidence in them once accepted. Together, these imply that block decoding wastes computation in a predictable way: as a block nears convergence, each full-width forward pass re-processes a stable accepted prefix in order to resolve a few remaining positions, work that could instead resolve tokens of the next block. We quantify this waste directly.

% We then introduce \algname (Section~\ref{sec:method}): the two rules governing the sampler, the recasting of the temporal temperature schedule as a fixed spatial profile over the window, and the resulting steady-state behavior of the window's entropy. We further show that \algname admits an interpretation as approximate speculative decoding, in which the parallel acceptances of each step act as a draft against a sequential target, the autoregressive-order decomposition of the diffusion model itself, and the entropy budget bounds the probability that an exact verifier would reject the draft (Section~\ref{subsec:specdec}).

% Finally, we evaluate \algname against the EB sampler on a wide range of benchmarks (Section~\ref{sec:experiments}), demonstrating consistently higher tokens per forward pass at matched quality, with only marginal degradation in the worst case.